\documentclass[12pt]{article}

\usepackage{fontspec}
\setmainfont{Times New Roman}
\usepackage{microtype}
\usepackage[a4paper,margin=2.5cm]{geometry}
\usepackage{setspace}
\onehalfspacing
\usepackage{longtable,booktabs,array,calc,tabularx}
\IfFileExists{makecell.sty}{\usepackage{makecell}}{}
\usepackage{caption}
\captionsetup{font=small,labelfont=bf}
\newcounter{none}
\usepackage{fancyvrb}
\usepackage{graphicx}
\usepackage{xcolor}
\usepackage{hyperref}
\hypersetup{
  colorlinks=true,
  linkcolor=black!70,
  urlcolor=black!70,
  citecolor=black!70
}
\IfFileExists{xurl.sty}{\usepackage{xurl}}{}
\setlength{\parindent}{1.25cm}
\setlength{\parskip}{0pt}
\providecommand{\tightlist}{%
  \setlength{\itemsep}{0pt}\setlength{\parskip}{0pt}}
\newlength{\cslhangindent}
\setlength{\cslhangindent}{1.5em}
\newlength{\csllabelwidth}
\setlength{\csllabelwidth}{3em}
\newenvironment{CSLReferences}[2]
 {\begin{list}{}{%
  \setlength{\itemindent}{0pt}
  \setlength{\leftmargin}{0pt}
  \setlength{\parsep}{0pt}
  \ifodd #1
    \setlength{\leftmargin}{\cslhangindent}
    \setlength{\itemindent}{-\cslhangindent}
  \fi
  \setlength{\itemsep}{#2\baselineskip}}}
 {\end{list}}
\newcommand{\CSLBlock}[1]{\hfill\break#1\hfill\break}
\newcommand{\CSLLeftMargin}[1]{\parbox[t]{\csllabelwidth}{#1}}
\newcommand{\CSLRightInline}[1]{\parbox[t]{\linewidth - \csllabelwidth}{#1}\break}
\newcommand{\CSLIndent}[1]{\hspace{\cslhangindent}#1}
\newcommand{\citeproctext}{}
\renewcommand{\abstractname}{Resumo}
\renewcommand{\contentsname}{Sumário}
\sloppy

\title{Portal Nacional de Contratações Públicas e governo aberto: uma
análise exploratória das capitais do Sudeste\\\large Transparência
formal, dados abertos e fragmentação institucional}
\author{PREENCHER\_NOME}
\date{01/07/2026}

\begin{document}
\pagenumbering{gobble}
\maketitle

\begin{abstract}
Este relatório analisa pregões eletrônicos publicados no Portal Nacional
de Contratações Públicas (PNCP) entre 15/06/2025 e 15/06/2026, com foco
nas prefeituras das capitais do Sudeste. O método combina quatro etapas:
coleta automatizada dos registros, análise de completude (verificação
de quais campos vêm efetivamente preenchidos), medição da facilidade de
consumo dos dados por meios automatizados e comparação institucional
entre as capitais. Os resultados indicam boa disponibilidade dos campos
básicos, ressalvas no valor homologado (o preço final aprovado) e no
link para o sistema de origem, além de forte fragmentação de CNPJs em
São Paulo, onde grande parte das contratações é publicada por unidades
municipais distintas da prefeitura central.
\end{abstract}

\begin{center}
Universidade de São Paulo\\
Governo Aberto\\
Docente: Jorge Machado
\end{center}

\tableofcontents

\thispagestyle{empty}
\newpage
\pagenumbering{arabic}

\section{Introdução}\label{introduuxe7uxe3o}

Este trabalho analisa o Portal Nacional de Contratações Públicas (PNCP)
como infraestrutura de governo aberto nas capitais do Sudeste
brasileiro. O PNCP é o sistema oficial em que órgãos públicos de todo o
país divulgam suas licitações e contratos; ele é especialmente relevante
porque a Lei nº 14.133/2021 (a nova Lei de Licitações) o institui como
ambiente nacional obrigatório de divulgação das contratações públicas
(Brasil, 2021).

A investigação foi organizada em duas questões de pesquisa
complementares. A primeira examina \emph{o que} o portal disponibiliza;
a segunda, \emph{como} esses dados podem ser obtidos:

Q1. Há completude nos dados fornecidos pelo PNCP nas prefeituras das
capitais dos estados do Sudeste brasileiro (São Paulo, Rio de Janeiro,
Belo Horizonte e Vitória)? Em outras palavras, os registros chegam ao
cidadão com seus campos efetivamente preenchidos?

Q2. Os dados das \textbf{APIs}\footnote{\textbf{API} (\emph{Application
  Programming Interface}, ou interface de programação de aplicações): é
  o canal pelo qual um programa consulta outro de forma automatizada,
  sem intervenção manual. Neste trabalho, é o meio de consultar os dados
  do PNCP diretamente por código.} do PNCP são facilmente consumíveis?
Isto é, um pesquisador ou jornalista consegue extrair esses dados em
escala sem obstáculos técnicos relevantes?

O argumento central que articula as duas questões é o seguinte: o PNCP
amplia a transparência \emph{formal} ao centralizar metadados (os dados
que descrevem cada contratação, como objeto, valores e datas),
documentos e links de origem; porém, sua efetividade como instrumento de
governo aberto depende de dois fatores adicionais --- a completude dos
registros (Q1) e a possibilidade prática de reconstituir a organização
administrativa de cada prefeitura a partir dos dados (Q2). O caso de São
Paulo evidencia essa tensão com mais nitidez: grande parte dos registros
municipais aparece fora do \textbf{CNPJ matriz}\footnote{\textbf{CNPJ
  matriz}: o CNPJ principal da prefeitura, que identifica o município
  como pessoa jurídica. Secretarias, fundos e empresas municipais podem
  ter CNPJs próprios (filiais ou entidades vinculadas) e publicar
  contratações em separado, fora do CNPJ matriz.} do município, de modo
que uma consulta restrita a esse CNPJ recuperaria apenas uma fração do
que a cidade efetivamente contratou.

\section{Objetivos e método}\label{objetivos-e-muxe9todo}

O objetivo geral é avaliar, de forma exploratória, como o PNCP expressa
princípios de governo aberto nas contratações das capitais do Sudeste. A
pesquisa combina análise documental e computacional, organizada em uma
sequência de etapas automatizadas (um \textbf{pipeline}\footnote{\textbf{Pipeline}:
  sequência encadeada de etapas automatizadas em que a saída de uma
  etapa alimenta a seguinte --- aqui, da coleta dos dados ao cálculo das
  métricas e à geração do relatório.}): coleta os registros por
\textbf{API}, transforma os dados em \textbf{snapshots}\footnote{\textbf{Snapshot}:
  cópia congelada dos dados no momento da coleta, salva em arquivo. Por
  ser fixa, permite reexecutar a análise e auditar os resultados sem
  depender de novas consultas à API, cujo conteúdo pode mudar com o
  tempo.} auditáveis, calcula métricas e gera relatórios reprodutíveis.

O recorte empírico considera pregões eletrônicos --- identificados como
modalidade 6 no PNCP --- publicados entre 15/06/2025 e 15/06/2026. Para
Rio de Janeiro, Belo Horizonte e Vitória, foram usados os CNPJs matriz
de cada município. Para São Paulo, combinou-se o CNPJ matriz com uma
varredura por UF (unidade da federação, ou seja, o estado) e por código
IBGE (o código numérico que o Instituto Brasileiro de Geografia e
Estatística atribui a cada município), filtrando apenas as entidades do
poder executivo municipal. Essa estratégia mais ampla para São Paulo foi
necessária justamente porque, como se verá, a cidade publica grande
parte de suas contratações fora do CNPJ matriz.

{\def\LTcaptype{none} % do not increment counter
\begin{longtable}[]{@{}llllll@{}}
\toprule\noalign{}
Capital & Fonte & Registros & Páginas & Total & Limite \\
\midrule\noalign{}
\endhead
\bottomrule\noalign{}
\endlastfoot
Sao Paulo & matriz & 30 & 12 & 12 & \\
Sao Paulo & município & 17654 & 359 & 359 & todas \\
Rio de Janeiro & matriz & 896 & 23 & 23 & \\
Belo Horizonte & matriz & 291 & 12 & 12 & \\
Vitoria & matriz & 234 & 12 & 12 & \\
\end{longtable}
}

A amostra principal inclui todos os registros elegíveis no período, após
\textbf{deduplicação}\footnote{\textbf{Deduplicação}: remoção de
  registros repetidos. Aqui, usou-se o campo
  \texttt{numeroControlePNCP}, o identificador único de cada contratação
  no portal, para garantir que cada pregão fosse contado uma só vez.}
pelo campo \texttt{numeroControlePNCP}. Para a análise de documentos,
foi usada uma subamostra determinística de até 100 registros por
capital. A subamostra é \emph{determinística} porque parte de uma
\textbf{semente}\footnote{\textbf{Semente} (\emph{seed}): número que
  inicializa o sorteio aleatório. Fixá-la (aqui, 20260608) faz com que a
  mesma amostra seja sorteada a cada execução, tornando o resultado
  reprodutível por terceiros.} fixa (20260608), o que permite reproduzir
exatamente o mesmo sorteio.

As métricas de Q1 medem a presença de campos essenciais --- objeto,
valores, datas, unidade administrativa, link de origem e documentos ---,
respondendo à pergunta sobre completude. As métricas de Q2 registram a
duração do experimento, o tempo médio de resposta e as falhas observadas
durante o consumo da \textbf{API}, respondendo à pergunta sobre
facilidade de consumo.

A reprodutibilidade foi organizada no repositório GitHub
\url{https://github.com/Amorim33/analise-pncp}. O repositório versiona
código Python, configurações, snapshots brutos, tabelas processadas,
métricas, análise exploratória e o relatório final em Markdown, LaTeX e
PDF, na branch \texttt{main}.

O processo foi assistido pelo \textbf{Codex}\footnote{\textbf{Codex}:
  agente de programação e documentação usado para implementar, validar e
  revisar o pipeline.} em etapas de programação e documentação, sob
supervisão humana. A divisão agêntica usou as \textbf{skills}\footnote{\textbf{Skills}:
  instruções locais que especializam agentes para tarefas como mapear a
  \textbf{API}, coletar dados, revisar amostragem e redigir o relatório.}
locais em \texttt{.agents/skills/}: mapeamento da \textbf{API}, coleta
de dados, metodologia de amostragem e redação do relatório. As decisões
substantivas, a validação das fontes e a interpretação final permanecem
sob responsabilidade do autor; a declaração de uso de IA está no
Apêndice B.

\section{Referencial teórico}\label{referencial-teuxf3rico}

Governo aberto costuma ser descrito como um arranjo que combina
transparência, participação social, colaboração e prestação de contas. A
\emph{Open Government Partnership} (OGP, ou Parceria para Governo
Aberto), iniciativa multilateral que reúne dezenas de países em torno do
tema, associa-o a compromissos de abertura, integridade pública e
participação cidadã (Open Government Partnership, 2011). De modo
convergente, a OCDE (Organização para a Cooperação e Desenvolvimento
Econômico) define governo aberto como uma cultura de governança que
promove a transparência, a integridade, a
\emph{accountability}\footnote{\emph{Accountability}: termo inglês sem
  tradução exata em português, que reúne as ideias de prestação de
  contas e de responsabilização --- o dever de quem exerce poder
  público de justificar seus atos e responder por eles.} e a
participação das partes interessadas (Organisation for Economic
Co-operation and Development, 2017).

Os dados abertos são a dimensão operacional desse debate --- isto é, o
mecanismo concreto pelo qual a transparência se realiza. Para que a
publicação seja útil, não basta disponibilizar informação em páginas
isoladas: é preciso permitir acesso, reuso, padronização e combinação
com outras bases (W3C Brasil, 2011)\footnote{O W3C (\emph{World Wide Web
  Consortium}) é o consórcio internacional que define os padrões da
  web.}. A Lei de Acesso à Informação reforça esse princípio ao tornar a
publicidade a regra e o sigilo a exceção (Brasil, 2011). Daí decorre o
critério que orienta este trabalho: a transparência pública deve ser
avaliada não apenas pela \emph{existência} dos dados, mas por sua
\emph{usabilidade} para fiscalização e pesquisa --- distinção que
justamente separa as duas questões de pesquisa (Q1 e Q2).

No caso brasileiro, o PNCP é uma infraestrutura particularmente
adequada a esse tipo de análise, pois centraliza informações e
documentos de contratações públicas (Brasil, 2026c). A existência de uma
API e de dados consultáveis amplia o potencial de reutilização por
cidadãos, pesquisadores e jornalistas (Brasil, 2026b, 2026a) --- desde
que, como argumentamos, esses dados se mostrem completos e efetivamente
acessíveis na prática.

\section{Desenvolvimento e análise}\label{desenvolvimento-e-anuxe1lise}

\subsection{Visão comparativa}\label{visuxe3o-comparativa}

{\def\LTcaptype{none} % do not increment counter
\begin{longtable}[]{@{}
  >{\raggedright\arraybackslash}p{(\linewidth - 10\tabcolsep) * \real{0.1667}}
  >{\raggedright\arraybackslash}p{(\linewidth - 10\tabcolsep) * \real{0.1667}}
  >{\raggedright\arraybackslash}p{(\linewidth - 10\tabcolsep) * \real{0.1667}}
  >{\raggedright\arraybackslash}p{(\linewidth - 10\tabcolsep) * \real{0.1667}}
  >{\raggedright\arraybackslash}p{(\linewidth - 10\tabcolsep) * \real{0.1667}}
  >{\raggedright\arraybackslash}p{(\linewidth - 10\tabcolsep) * \real{0.1667}}@{}}
\toprule\noalign{}
\begin{minipage}[b]{\linewidth}\raggedright
Capital
\end{minipage} & \begin{minipage}[b]{\linewidth}\raggedright
Candidatos
\end{minipage} & \begin{minipage}[b]{\linewidth}\raggedright
Amostra
\end{minipage} & \begin{minipage}[b]{\linewidth}\raggedright
CNPJs distintos
\end{minipage} & \begin{minipage}[b]{\linewidth}\raggedright
Registros na matriz
\end{minipage} & \begin{minipage}[b]{\linewidth}\raggedright
Participação da matriz
\end{minipage} \\
\midrule\noalign{}
\endhead
\bottomrule\noalign{}
\endlastfoot
Sao Paulo & 2662 & 2662 & 67 & 30 & 1.1\% \\
Rio de Janeiro & 896 & 896 & 1 & 896 & 100.0\% \\
Belo Horizonte & 291 & 291 & 1 & 291 & 100.0\% \\
Vitoria & 234 & 234 & 1 & 234 & 100.0\% \\
\end{longtable}
}

A tabela acima já revela a primeira diferença relevante: a concentração
institucional, medida pela coluna ``Participação da matriz'' (a fração
dos registros publicada sob o CNPJ matriz). Rio de Janeiro, Belo
Horizonte e Vitória concentram 100\% dos registros no CNPJ matriz no
recorte adotado --- ou seja, consultar o CNPJ da prefeitura recupera
toda a atividade contratual. São Paulo, ao contrário, apresenta apenas
1,1\% sob o CNPJ matriz e 67 CNPJs distintos, associados a unidades como
fundo municipal, hospital, companhia de engenharia de tráfego,
secretaria e subprefeitura. É essa dispersão que a próxima subseção
detalha.

\subsection{Fragmentação em São
Paulo}\label{fragmentauxe7uxe3o-em-suxe3o-paulo}

{\def\LTcaptype{none} % do not increment counter
\begin{longtable}[]{@{}ll@{}}
\toprule\noalign{}
Métrica & Valor \\
\midrule\noalign{}
\endhead
\bottomrule\noalign{}
\endlastfoot
CNPJ matriz & 46395000000139 \\
Registros candidatos & 2662 \\
Registros no CNPJ matriz & 30 \\
Registros fora do CNPJ matriz & 2632 \\
Participação fora da matriz & 98.9\% \\
CNPJs distintos & 67 \\
CNPJs fora da matriz & 66 \\
\end{longtable}
}

Os números são expressivos: dos 2.662 registros candidatos, apenas 30
(1,1\%) estão sob o CNPJ matriz e 2.632 (98,9\%) estão distribuídos por
outros 66 CNPJs municipais. Esse resultado é substantivo para o governo
aberto. Uma consulta limitada ao CNPJ matriz de São Paulo recuperaria
menos de 2\% das contratações elegíveis e, portanto, daria ao cidadão
uma imagem gravemente subestimada da atividade do município. A
transparência formal do PNCP convive, assim, com uma barreira de
reconstrução institucional: para enxergar o todo, o cidadão precisa
saber de antemão que a prefeitura pode publicar contratações por
dezenas de CNPJs distintos --- conhecimento que o próprio portal não lhe
fornece de forma direta.

\subsection{Q1: completude dos dados}\label{q1-completude-dos-dados}

{\def\LTcaptype{none} % do not increment counter
\begin{longtable}[]{@{}
  >{\raggedright\arraybackslash}p{(\linewidth - 14\tabcolsep) * \real{0.1250}}
  >{\raggedright\arraybackslash}p{(\linewidth - 14\tabcolsep) * \real{0.1250}}
  >{\raggedright\arraybackslash}p{(\linewidth - 14\tabcolsep) * \real{0.1250}}
  >{\raggedright\arraybackslash}p{(\linewidth - 14\tabcolsep) * \real{0.1250}}
  >{\raggedright\arraybackslash}p{(\linewidth - 14\tabcolsep) * \real{0.1250}}
  >{\raggedright\arraybackslash}p{(\linewidth - 14\tabcolsep) * \real{0.1250}}
  >{\raggedright\arraybackslash}p{(\linewidth - 14\tabcolsep) * \real{0.1250}}
  >{\raggedright\arraybackslash}p{(\linewidth - 14\tabcolsep) * \real{0.1250}}@{}}
\toprule\noalign{}
\begin{minipage}[b]{\linewidth}\raggedright
Capital
\end{minipage} & \begin{minipage}[b]{\linewidth}\raggedright
Obj.
\end{minipage} & \begin{minipage}[b]{\linewidth}\raggedright
Estim.
\end{minipage} & \begin{minipage}[b]{\linewidth}\raggedright
Homol.
\end{minipage} & \begin{minipage}[b]{\linewidth}\raggedright
Data
\end{minipage} & \begin{minipage}[b]{\linewidth}\raggedright
Unid.
\end{minipage} & \begin{minipage}[b]{\linewidth}\raggedright
Link
\end{minipage} & \begin{minipage}[b]{\linewidth}\raggedright
Docs
\end{minipage} \\
\midrule\noalign{}
\endhead
\bottomrule\noalign{}
\endlastfoot
Sao Paulo & 100.0\% & 100.0\% & 60.4\% & 100.0\% & 100.0\% & 99.8\% &
100.0\% \\
Rio de Janeiro & 100.0\% & 100.0\% & 62.8\% & 100.0\% & 100.0\% &
100.0\% & 100.0\% \\
Belo Horizonte & 100.0\% & 100.0\% & 78.0\% & 100.0\% & 100.0\% & 96.9\%
& 100.0\% \\
Vitoria & 100.0\% & 100.0\% & 66.7\% & 100.0\% & 100.0\% & 0.0\% &
100.0\% \\
\end{longtable}
}

A tabela apresenta, para cada capital, o percentual de registros em que
cada campo veio preenchido: objeto (Obj.), valor estimado (Estim.),
valor homologado (Homol.), data, unidade administrativa (Unid.), link de
origem (Link) e documentos (Docs). A resposta a Q1 é positiva para os
campos estruturantes, mas com ressalvas relevantes. O objeto, o valor
estimado, as datas básicas e a unidade administrativa apareceram em
praticamente 100\% dos registros elegíveis. A variação concentrou-se em
três campos: o \textbf{valor homologado} (o preço final aprovado ao fim
do pregão), o \textbf{link do sistema de origem} (o endereço que remete
ao sistema próprio do município, de onde o registro foi enviado ao PNCP)
e os documentos. Essa diferença importa para o controle social porque,
sem o valor homologado, não se sabe quanto foi efetivamente contratado;
e, sem o link de origem, perde-se o caminho para auditar o processo em
sua fonte. Vitória ilustra o ponto extremo: 0\% de link de origem,
apesar de 100\% de documentos anexados.

Exemplos pseudoaleatórios\footnote{``Pseudoaleatórios'' porque sorteados
  por um gerador de números pseudoaleatórios a partir da semente fixa
  --- o sorteio parece aleatório, mas é reproduzível.} da subamostra
documental ilustram essa variação:

\begin{itemize}
\tightlist
\item
  Sao Paulo (\texttt{13864377000130-1-002050/2025}): data 18/11/2025;
  valor homologado presente; link de origem presente; 1 documento(s);
  objeto: REGISTRO De PREÇOS PARA O FORNECIMENTO CANULA TRAQUEOSTOMIA
  COM BALAO, DESCARTAVEL, ESTERIL - NR. 4,5 E SONDA ENDOTRAQUEAL, COM
  BALAO, RADIOPACO, D\ldots{}
\item
  Rio de Janeiro (\texttt{42498733000148-1-001315/2025}): data
  24/07/2025; valor homologado presente; link de origem presente; 1
  documento(s); objeto: Aquisição de Equipamento de Proteção Individual
  --EPI
\item
  Belo Horizonte (\texttt{18715383000140-1-000189/2026}): data
  27/03/2026; valor homologado presente; link de origem presente; 1
  documento(s); objeto: O objeto da presente licitação é o registro de
  preços para aquisição de materiais de higiene e limpeza específicos,
  conforme condições e exigências\ldots{}
\item
  Vitoria (\texttt{27142058000126-1-000676/2025}): data 28/11/2025;
  valor homologado presente; link de origem ausente; 6 documento(s);
  objeto: REGISTRO DE PREÇOS VISANDO FUTURO E EVENTUAL FORNECIMENTO DE
  ÁGUA MINERAL GALÃO 20L
\end{itemize}

{\def\LTcaptype{none} % do not increment counter
\begin{longtable}[]{@{}llllll@{}}
\toprule\noalign{}
Capital & Com documentos & Amostra & Mín. & Máx. & Média \\
\midrule\noalign{}
\endhead
\bottomrule\noalign{}
\endlastfoot
Sao Paulo & 100 & 100 & 1 & 1 & 1.0 \\
Rio de Janeiro & 100 & 100 & 1 & 1 & 1.0 \\
Belo Horizonte & 100 & 100 & 1 & 2 & 1.0 \\
Vitoria & 100 & 100 & 1 & 12 & 5.3 \\
\end{longtable}
}

A tabela acima detalha quantos documentos cada registro traz anexados na
subamostra. Em todas as capitais, todos os 100 itens da subamostra
tinham ao menos um documento. A diferença está na quantidade: São Paulo
e Rio de Janeiro anexam, em regra, um único documento por registro,
enquanto Vitória apresenta média de 5,3 e chega a 12 documentos em um
mesmo registro. Mais documentos tendem a significar mais contexto
disponível ao cidadão, embora a leitura conjunta com a tabela anterior
mostre que volume de documentos não compensa a ausência do link de
origem.

\subsection{Q2: consumo da API}\label{q2-consumo-da-api}

{\def\LTcaptype{none} % do not increment counter
\begin{longtable}[]{@{}llllll@{}}
\toprule\noalign{}
Etapa & Req. & Sucesso & Média & Duração & Falhas \\
\midrule\noalign{}
\endhead
\bottomrule\noalign{}
\endlastfoot
Snapshots & n/a & n/a & n/a & n/a & n/a \\
Tentativa live & 6 & 0 & n/a & 2min 56.6s & 6 \\
Documentos & 400 & 400 & 0.10s & 40.51s & 0 \\
Total & & & & 3min 37.8s & \\
\end{longtable}
}

A resposta a Q2 é intermediária. Por um lado, a API é consumível por
programas: retorna os dados em \textbf{JSON}\footnote{\textbf{JSON}
  (\emph{JavaScript Object Notation}): formato de texto estruturado,
  legível por máquina e por humano, amplamente usado para troca de dados
  entre sistemas.} estruturado e oferece \textbf{paginação}\footnote{\textbf{Paginação}:
  entrega dos resultados em ``páginas'' sucessivas em vez de tudo de uma
  vez, exigindo várias requisições para obter o conjunto completo.}. Por
outro, o consumo não é trivial em uma janela anual: foi necessário
dividir as consultas em blocos mensais, inserir pausas entre as
requisições e tratar o \emph{limite de taxa} (\emph{rate limit}) ---
restrição que o servidor impõe ao número de requisições aceitas em um
dado intervalo. A duração total do experimento, registrada na tabela
acima, soma a coleta das contratações, a geração da amostra e a consulta
aos documentos da subamostra.

Quanto à estabilidade, os snapshots reutilizados não registraram
nenhuma falha persistente na coleta bem-sucedida anterior, e a consulta
de documentos da execução final também terminou sem falhas. A ressalva
está na linha ``Tentativa live'' da tabela: a tentativa de coletar os
dados ao vivo nesta execução falhou nas seis requisições, e o pipeline
recorreu, por segurança, aos snapshots já existentes. O servidor
respondeu com o código \textbf{HTTP 503} (\emph{Service Unavailable},
serviço indisponível), indicando indisponibilidade temporária do lado do
PNCP --- e não um erro do código de coleta.

Ao longo da experimentação, três obstáculos técnicos recorrentes foram
identificados e contornados:

\begin{itemize}
\tightlist
\item
  \textbf{HTTP 429} (\emph{Too Many Requests}): o servidor recusava
  chamadas repetidas por excesso de requisições. Tratado com
  \emph{backoff} (espera progressiva entre tentativas) e nova
  tentativa.
\item
  \emph{Timeouts} (estouros do tempo limite de resposta) na paginação de
  janelas anuais longas. Mitigados pela coleta em blocos
  (\emph{chunks}) mensais de até 31 dias.
\item
  Risco de o servidor devolver uma página HTML com código \textbf{HTTP
  200} (que normalmente sinaliza sucesso), quando o esperado seria JSON.
  Tratado pela validação do \emph{Content-Type} (o tipo de conteúdo
  declarado pela resposta) antes de interpretar o dado como JSON.
\end{itemize}

\subsection{Constatações
empíricas}\label{constatauxe7uxf5es-empuxedricas}

\begin{itemize}
\tightlist
\item
  Na amostra candidata de São Paulo, 98,9\% dos registros elegíveis
  ficaram fora do CNPJ matriz; nas demais capitais analisadas, os
  registros por CNPJ matriz concentraram 100\% dos candidatos coletados.
\item
  São Paulo apresentou 67 CNPJs distintos no recorte elegível, enquanto
  Rio de Janeiro, Belo Horizonte e Vitória apareceram com um CNPJ cada
  no recorte por matriz.
\item
  Vitória teve documentos vinculados em todos os itens da subamostra
  documental, mas nenhum dos registros elegíveis trouxe o campo
  \texttt{linkSistemaOrigem} (o link para o sistema de origem), o que
  reduz a rastreabilidade até a fonte.
\item
  O valor homologado apareceu com menor completude em São Paulo
  (60,4\%), indicando que parte dos registros estava em fase anterior ao
  resultado ou sem homologação registrada no recorte.
\item
  A quantidade de documentos anexados variou de forma relevante: em
  Vitória, um item da amostra chegou a 12 documentos, enquanto os demais
  municípios tiveram padrão mais concentrado (em geral, um documento por
  registro).
\end{itemize}

\section{Conclusões}\label{conclusuxf5es}

Retomando as duas questões de pesquisa: quanto à completude (Q1), a
análise indica que o PNCP é uma infraestrutura importante de governo
aberto, pois centraliza registros, expõe campos padronizados com alta
disponibilidade (objeto, valores estimados, datas e unidade), oferece
documentos e permite automação por API. Quanto à facilidade de consumo
(Q2), a API é utilizável, mas exige cuidados técnicos --- divisão em
blocos, controle de taxa e tratamento de falhas --- que elevam a barreira
para o usuário menos especializado. Em conjunto, esses elementos
fortalecem a transparência formal e criam condições para o controle
social, ainda que de forma desigual.

Ao mesmo tempo, a comparação mostra que abertura não é sinônimo
automático de inteligibilidade. A fragmentação de CNPJs em São Paulo
torna a busca mais complexa e pode subestimar drasticamente a atividade
contratual municipal quando o pesquisador consulta apenas o CNPJ matriz
--- como demonstram os 98,9\% de registros publicados fora dele. A
principal contribuição do trabalho está, portanto, em evidenciar que a
avaliação de governo aberto deve observar não só a presença e a
acessibilidade dos dados, mas também a sua arquitetura institucional:
quem publica o quê, e sob qual identidade.

Como conclusão regional exploratória, as capitais do Sudeste analisadas
apresentam boa disponibilidade básica de dados e documentos no PNCP, mas
diferem na forma de organização dos registros. A agenda de melhoria
passa por documentação mais clara dos CNPJs/unidades, maior completude
de links de origem e mecanismos que facilitem ao cidadão reconstruir o
universo de órgãos vinculados a cada prefeitura.

\section{Referências}\label{referuxeancias}

\protect\phantomsection\label{refs}
\begin{CSLReferences}{1}{0}
\bibitem[\citeproctext]{ref-lei12527}
Brasil. (2011). \emph{Lei nº 12.527, de 18 de novembro de 2011}.
\url{https://www.planalto.gov.br/ccivil_03/_ato2011-2014/2011/lei/l12527.htm}

\bibitem[\citeproctext]{ref-lei14133}
Brasil. (2021). \emph{Lei nº 14.133, de 1º de abril de 2021}.
\url{https://www.planalto.gov.br/ccivil_03/_ato2019-2022/2021/lei/l14133.htm}

\bibitem[\citeproctext]{ref-pncpApiConsulta}
Brasil. (2026a). \emph{API PNCP Consulta}.
\url{https://pncp.gov.br/pncp-consulta/v3/api-docs}

\bibitem[\citeproctext]{ref-pncpDadosAbertos}
Brasil. (2026b). \emph{PNCP: dados abertos}.
\url{https://www.gov.br/pncp/pt-br/acesso-a-informacao/dados-abertos}

\bibitem[\citeproctext]{ref-pncpPortal}
Brasil. (2026c). \emph{Portal Nacional de Contratações Públicas}.
\url{https://www.gov.br/pncp/pt-br/pncp}

\bibitem[\citeproctext]{ref-ogpDeclaration}
Open Government Partnership. (2011). \emph{Open Government Declaration}.
\url{https://www.opengovpartnership.org/process/joining-ogp/open-government-declaration/}

\bibitem[\citeproctext]{ref-oecd2017}
Organisation for Economic Co-operation and Development. (2017).
\emph{Recommendation of the Council on Open Government}.
\url{https://legalinstruments.oecd.org/en/instruments/OECD-LEGAL-0438}

\bibitem[\citeproctext]{ref-w3cManualDadosAbertos}
W3C Brasil. (2011). \emph{Manual dos dados abertos: governo}.
\url{https://www.w3c.br/pub/Materiais/PublicacoesW3C/Manual_Dados_Abertos_WEB.pdf}

\end{CSLReferences}

\appendix

\section{Apêndice A: exemplos compactos da
API}\label{apuxeandice-a-exemplos-compactos-da-api}

\subsection{Sao Paulo - exemplo fora do CNPJ
matriz}\label{sao-paulo---exemplo-fora-do-cnpj-matriz}

\begin{itemize}
\tightlist
\item
  Controle PNCP: \texttt{01164292000160-1-000020/2026}
\item
  CNPJ/órgão: \texttt{01164292000160} --- MUNICIPIO DE CACU
\item
  Unidade: PREFEITURA MUNICIPAL DE CAÇU
\item
  Valor estimado: 137600.4
\item
  Valor homologado: não informado
\item
  Situação: Divulgada no PNCP
\item
  Documentos listados no exemplo: 0
\item
  Objeto: CONTRATAÇÃO DE EMPRESA DO RAMO PARA A AQUISIÇÃO E INSTALAÇÃO
  DE PARQUE INFANTIL, DEVIDAMENTE CERTIFICADO POR ÓRGÃO COMPETENTE, PARA
  ATENDER AS NECESSIDADES DO FUNDO MUNICIPAL DE EDUCAÇÃO/FME, CONFORME
  EMENDA PARLAMENTAR Nº 1299.2/2025, PROCESSO Nº 202500005\ldots{}
\end{itemize}

\subsection{Sao Paulo - exemplo no CNPJ
matriz}\label{sao-paulo---exemplo-no-cnpj-matriz}

\begin{itemize}
\tightlist
\item
  Controle PNCP: \texttt{46395000000139-1-000007/2026}
\item
  CNPJ/órgão: \texttt{46395000000139} --- MUNICIPIO DE SAO PAULO
\item
  Unidade: PMSP - SECRETARIA DO GOVERNO MUNICIPAL
\item
  Valor estimado: 62640.7
\item
  Valor homologado: 42066.57
\item
  Situação: Divulgada no PNCP
\item
  Documentos listados no exemplo: 0
\item
  Objeto: Aquisição de materiais de construção, abrangendo itens das
  linhas hidráulica, elétrica, ferragens, pintura, alvenaria,
  acabamentos e correlatos, conforme as especificações, quantidades e
  condições estabelecidas neste Termo de Referência considerando
  AUSÊNCI\ldots{}
\end{itemize}

\subsection{Rio de Janeiro}\label{rio-de-janeiro}

\begin{itemize}
\tightlist
\item
  Controle PNCP: \texttt{42498733000148-1-000104/2026}
\item
  CNPJ/órgão: \texttt{42498733000148} --- MUNICIPIO DE RIO DE JANEIRO
\item
  Unidade: PREFEITURA MUNICIPAL DO RIO DE JANEIRO - RJ
\item
  Valor estimado: 0.0
\item
  Valor homologado: não informado
\item
  Situação: Divulgada no PNCP
\item
  Documentos listados no exemplo: 0
\item
  Objeto: CoNTRATAÇÃO DE EMPRESA ESPECIALIZADA PARA PRESTAÇÃO DE
  SERVIÇOS DE PLANEJAMENTO, ORGANIZAÇÃO, EXECUÇÃO, ACOMPANHAMENTO,
  FORNECIMENTO DE BENS E INSUMOS, INFRAESTRUTURA E APOIO LOGÍSTICO,
  VISANDO A REALIZAÇÃO DO BAILE DA CINELÂNDIA E DESFILES DA AVENIDA
  CHILE\ldots{}
\end{itemize}

\subsection{Belo Horizonte}\label{belo-horizonte}

\begin{itemize}
\tightlist
\item
  Controle PNCP: \texttt{18715383000140-1-000024/2026}
\item
  CNPJ/órgão: \texttt{18715383000140} --- MUNICIPIO DE BELO HORIZONTE
\item
  Unidade: PREF.MUN.DE BELO HORIZONTE
\item
  Valor estimado: 3585413.35
\item
  Valor homologado: 1969468.81
\item
  Situação: Divulgada no PNCP
\item
  Documentos listados no exemplo: 0
\item
  Objeto: O objeto da presente licitação é a aquisição de COLETES DE
  PROTEÇÃO BALÍSTICA FLEXÍVEL, PARA USO POLICIAL, NÍVEL DE PROTEÇÃO
  III-A, COM CAPA EXTERNA TIPO MODULAR (COM CAPA SOBRESSALENTE TIPO
  MODULAR), com a finalidade de atender atividades operacionais da
  G\ldots{}
\end{itemize}

\section{Apêndice B: declaração de uso de
IA}\label{apuxeandice-b-declarauxe7uxe3o-de-uso-de-ia}

Este trabalho utilizou ferramentas de inteligência artificial generativa
para apoio à programação do pipeline de coleta e análise, organização
dos resultados, revisão textual e geração inicial do relatório. A
seleção do tema, a validação das fontes, a interpretação dos resultados
e a versão final são de responsabilidade do(s) autor(es).

\section{Apêndice C: limitações
metodológicas}\label{apuxeandice-c-limitauxe7uxf5es-metodoluxf3gicas}

A pesquisa é exploratória e não representa todos os municípios do
Sudeste. A API de consulta por publicação limita cada requisição a
janelas de até 365 dias; por isso, o recorte usa a maior janela aceita
em uma consulta reprodutível. Além disso, documentos vinculados foram
analisados por subamostra determinística, e os resultados podem mudar
com novas publicações, retificações ou alterações na API do PNCP.

\end{document}
